% ==================================================
%  Resumen
% ==================================================
\noindent\textbf{Resumen.}

Este trabajo desarrolla y evalúa un prototipo heurístico de diseño inverso asistido por aprendizaje automático para metalentes plasmónicas de nano-rendijas, con el objetivo de guiar de manera eficiente la búsqueda de geometrías candidatas que enfoquen luz a $\lambda = 630$~nm, reduciendo el número de evaluaciones requeridas mediante simulación electromagnética directa (FDTD). A partir de un conjunto de $1{,}110$ simulaciones FDTD, se entrenó un modelo subrogado que, tras un re-encuadre metodológico, predice el campo óptico focal complejo (amplitud y fase) en lugar de un escalar de calidad. Esta reformulación permitió aprovechar información densa de todas las simulaciones y evitar el desbalance estructural que limitó el enfoque inicial basado en la predicción directa de $R^2$.

El modelo subrogado se empleó como evaluador rápido y criterio de priorización dentro de un algoritmo genético restringido al vecindario de geometrías previamente validadas, complementado con un refinamiento final mediante simulación FDTD directa. En este sentido, el sistema no se plantea como un predictor global de geometrías óptimas en todo el espacio de diseño, sino como un pipeline de búsqueda local asistida por aprendizaje automático, útil para explorar y refinar regiones prometedoras con menor costo computacional. Bajo este esquema, se obtuvo una geometría recomendada (\texttt{hc627\_036}) que cumplió el criterio operativo de éxito, definido como $R^2 \geq 0.90$ en el ajuste parabólico del perfil de intensidad y confirmado por FDTD, alcanzando $R^2_{\text{real}} = 0.9256$.

El análisis posterior mostró que el paisaje de diseño contiene cuencas de desempeño cualitativamente distintas: algunas de tipo ``aguja'', con máximos altos pero sensibles a pequeñas perturbaciones geométricas, y otras de tipo más ``amplio'', con desempeño relativamente más estable ante perturbaciones gaussianas de escala nanométrica. En particular, la región asociada a la geometría recomendada mostró mayor robustez relativa que otras soluciones de alto $R^2$, lo que la hace más atractiva desde una perspectiva de ingeniería. Finalmente, durante la fase de búsqueda, el surrogate permitió evaluar geometrías seis órdenes de magnitud más rápido que una simulación FDTD individual, confirmando su utilidad como herramienta de aceleración y priorización, aunque no como sustituto definitivo de la validación electromagnética directa.


\vspace{0.5em}
\noindent\textbf{Palabras clave:} metalentes plasmónicas, diseño inverso, modelos subrogados, búsqueda local asistida, algoritmo genético, FDTD.
